\section{Morfismo de incidencia hexada--octada y conservación de los canales orientados \(90/120\)}

Fijados la dodecada y el alineamiento del
\cref{sec:w12-w24-elevacion}, sea \(H\) una hexada en las coordenadas HMT,
sea \(\ell(H)\subset D\) su imagen y sea
\(O_H=\iota_D(\ell(H))\) su elevación única. Al marcar un punto
y un par de puntos de la hexada residual se obtienen
\[
F_6(H)=\ell(H)\times\binom{\ell(H)}2,
\qquad
|F_6|=6\binom62=90.
\]
Si el punto marcado puede recorrer la octada:
\[
F_8(H)=O_H\times\binom{\ell(H)}2,
\qquad
|F_8|=8\binom62=120.
\]
La inclusión \(\ell(H)\hookrightarrow O_H\) induce el morfismo de incidencia
\[
 \mathcal I_{6\to8}:F_6(H)\hookrightarrow F_8(H).
\]
Al normalizar la diferencia por el factor
elegido \(24\binom62\), se obtiene
\[
\boxed{
\frac{90-120}{24\binom62}
=\frac{6-8}{24}
=-\frac1{12}.}
\]
La igualdad es un defecto normalizado de la interfaz combinatoria
\(6\to8\). El factor \(24\binom62\) forma parte de la definición de esta
normalización; la inclusión por sí sola determina la diferencia \(-30\).
El operador \(\mathcal I_{6\to8}\) tiene como dominio configuraciones
marcadas de incidencia; no es el extractor transversal del estado
dodecafásico ni la familia de momentos angulares.
